\SourceLeaf{395}%
\[
\begin{array}{l}
M\\
M-2(a-r\cos A)\sin(A''-A')\\
M-2(a'-r\cos A')\sin(A-A'')\\
M-2(a''-r\cos A'')\sin(A'-A)
\end{array}
\]
We thus clearly have an equation of the form
\[
\gamma+\delta r+\varepsilon rr+\zeta r^3=\theta\theta-uu=4\alpha\alpha\beta\beta
\]
where \(\gamma,\delta,\varepsilon,\zeta\) are given quantities. The condition for a maximum then clearly determines \(r\) from the following quadratic equation
\[
\delta+2\varepsilon r+3\zeta rr=0
\]
\par
The coefficients of this equation are found still more easily by the following consideration. For any ellipse tangent to the four sides of the quadrilateral whose centre has abscissa \(r\), four times the product of the squares of its semimajor and semiminor axes is in general
\[
=\gamma+\delta r+\varepsilon rr+\zeta r^3
\]
and this expression must give \(=0\) when a value corresponding to one of the three vanishing ellipses considered above is substituted for \(r\). These three values are the distances from the origin of the abscissas to the two midpoints of the diagonals of the quadrilateral and to the midpoint of the straight segment joining the intersections of the two pairs of sides of the quadrilateral. I denote these three points by \(C,D,E\) and their abscissas by \(c,d,e\). Thus
\[
r^3+\tfrac{\varepsilon}{\zeta}rr+\tfrac{\delta}{\zeta}r+\tfrac{\gamma}{\zeta}
\]
must clearly be identical to the product \((r-c)(r-d)(r-e)\). Consequently the quadratic equation above is
\[
3rr-2r(c+d+e)+cd+ce+de=0
\]
Its roots are
\[
\tfrac{c+d+e}{3}+\tfrac13\BareRoot{(cc+dd+ee+cd+ce+de)}
\]
